% ---------------------------------------------------------------------------
% Minimal paper template — paper-skills
%
% Build:  pdflatex main → bibtex main → pdflatex main → pdflatex main
% Or, in VS Code with the "LaTeX Workshop" extension, just save the file.
%
% Bibliography: cite-keys come from refs.bib. The /bib-search skill (run with
% --for-paper <this-folder-name>) appends new entries to refs.bib for you.
% ---------------------------------------------------------------------------
\documentclass[11pt]{article}

\usepackage[utf8]{inputenc}
\usepackage[T1]{fontenc}
\usepackage{lmodern}
\usepackage[margin=1in]{geometry}
\usepackage{microtype}
\usepackage{hyperref}
\usepackage{graphicx}
\usepackage{booktabs}
\usepackage{natbib}            % \citep{} and \citet{}
\bibliographystyle{plainnat}

\title{Your Paper Title}
\author{Author One \and Author Two}
\date{\today}

\begin{document}
\maketitle

\begin{abstract}
One paragraph stating the problem, your contribution, and the result. Write this
last. Keep it self-contained --- a reader should understand the paper's claim
from the abstract alone.
\end{abstract}

\section{Introduction}
\label{sec:intro}
Motivate the problem and state your contribution. Cite related work with
\citep{example-key} for parenthetical citations and \citet{example-key} when the
author is the sentence subject.

\section{Related Work}
\label{sec:related}
Position your work against the literature. Use \verb|/bib-search| to discover
references and \verb|/claim-cite| to find the strongest citation for a specific
sentence.

\section{Method}
\label{sec:method}
Describe your approach.

\section{Results}
\label{sec:results}
Present the evidence.

\section{Conclusion}
\label{sec:conclusion}
Restate the contribution and name the limitations.

\bibliography{refs}

\end{document}
